\documentclass[aspectratio=169]{beamer}
\usetheme{Madrid}
\usecolortheme{default}
\usepackage{hyperref}
\usepackage{graphicx}
\usepackage{booktabs}

\title{The Data \& AI Challenge}
\subtitle{Rethinking Candidate Discovery: Beyond Keywords}
\author{Team: TOPGUN}
\date{}

\begin{document}

% Slide 1
\begin{frame}
  \titlepage
\end{frame}

% Slide 2 (Team Members)
\begin{frame}{Team TOPGUN}
  \begin{itemize}
    \item \textbf{Nethavath Praveen} (Leader) -- \href{mailto:praveeeening@gmail.com}{praveeeening@gmail.com}
    \item \textbf{Abhinaya Nunna} -- \href{mailto:abhinayaseven@gmail.com}{abhinayaseven@gmail.com}
    \item \textbf{Botsa Radesh} -- \href{mailto:radeshbotsa@gmail.com}{radeshbotsa@gmail.com}
    \item \textbf{Sai Sameer Killamsetti} -- \href{mailto:saisameerkillamsetty8@gmail.com}{saisameerkillamsetty8@gmail.com}
  \end{itemize}
\end{frame}

% Slide 2
\begin{frame}{Executive Summary}
  \begin{itemize}
    \item \textbf{Objective:} Build an intelligent candidate discovery and ranking system.
    \item \textbf{Core Philosophy:} Rank candidates based on semantic meaning, career trajectory, and behavioral intent, rather than basic keyword frequency.
    \item \textbf{Performance:} Process 100,000+ candidates in under 35 seconds.
    \item \textbf{Architecture:} Fully offline, deterministic, and highly scalable.
    \item \textbf{Transparency:} Integrated Explainable AI (XAI) for every scored profile.
  \end{itemize}
\end{frame}

% Slide 3
\begin{frame}{The Problem Statement}
  \framesubtitle{The Limitations of Traditional ATS}
  \begin{itemize}
    \item \textbf{Keyword Reliance:} Traditional systems count how many times a term appears.
    \item \textbf{Lack of Context:} "Used Python in college" vs. "Architected a Python microservice" are scored identically.
    \item \textbf{Recruiter Fatigue:} Teams waste hours sifting through falsely-ranked profiles.
    \item \textbf{Missed Talent:} Exceptional engineers who write concise resumes get buried beneath keyword-stuffed profiles.
  \end{itemize}
\end{frame}

% Slide 4
\begin{frame}{The Flaws of Current ATS Systems}
  \framesubtitle{"Keyword Stuffers" vs. True Experts}
  \begin{itemize}
    \item \textbf{The Stuffer:} Pastes the entire job description in white font or tiny text. Scores 100\% on legacy ATS.
    \item \textbf{The Expert:} Focuses on impact, system design, and leadership. Often omits basic technology tags they mastered a decade ago.
    \item \textbf{The Consequence:} A fundamental misalignment between the tool's output and the recruiter's actual need.
  \end{itemize}
\end{frame}

% Slide 5
\begin{frame}{Our Solution Overview}
  \begin{itemize}
    \item \textbf{Semantic Understanding:} We group skills by meaning, not string matches.
    \item \textbf{Career Shape Analysis:} We look at the trajectory of titles over time.
    \item \textbf{Intent Scoring:} We measure how actively the candidate is looking.
    \item \textbf{Anti-Gaming Integrity:} We statistically detect and penalize artificial inflation.
    \item \textbf{The Result:} A rank order that accurately mimics an expert human recruiter.
  \end{itemize}
\end{frame}

% Slide 6
\begin{frame}{The 4 Pillars of Intelligent Ranking}
  Our engine evaluates candidates across four heavily weighted axes:
  \begin{enumerate}
    \item \textbf{Semantic Competency (40\%):} Technical fit across clustered concepts.
    \item \textbf{Career Trajectory (30\%):} Pace of promotion and company pedigree.
    \item \textbf{Behavioral Intent (20\%):} Engagement signals and readiness to move.
    \item \textbf{Location Intelligence (10\%):} Geographical alignment with the role.
  \end{enumerate}
\end{frame}

% Slide 7
\begin{frame}{Pillar 1: Semantic Competency (40\%)}
  \framesubtitle{Moving Beyond Exact String Matches}
  \begin{itemize}
    \item \textbf{The Goal:} Recognize category expertise.
    \item \textbf{Mechanism:} If a job requires "Vector Databases," candidates get credit if they list Pinecone, Weaviate, FAISS, or ChromaDB.
    \item \textbf{Score Calculation:} Proportional to the number of critical semantic clusters a candidate possesses.
    \item \textbf{Bonus Allocation:} Additional weighting for full-stack capability across multiple disjoint clusters.
  \end{itemize}
\end{frame}

% Slide 8
\begin{frame}{Skill Clusters in Practice}
  \begin{table}
    \centering
    \begin{tabular}{@{}ll@{}}
      \toprule
      \textbf{Cluster Category} & \textbf{Accepted Technologies} \\
      \midrule
      Core Programming & Python, Go, C++, Java \\
      Vector Databases & Pinecone, Weaviate, Qdrant, Milvus, FAISS \\
      Retrieval \& Ranking & BM25, Cross-encoders, NDCG, Hybrid Search \\
      LLM Engineering & Fine-tuning, LoRA, LangChain, HuggingFace \\
      Infrastructure & Docker, Kubernetes, FastAPI, AWS, GCP \\
      \bottomrule
    \end{tabular}
  \end{table}
\end{frame}

% Slide 9
\begin{frame}{Pillar 2: Career Trajectory (30\%)}
  \framesubtitle{Promotion Velocity and Experience Bands}
  \begin{itemize}
    \item \textbf{Experience Banding:} Evaluates if the total years of experience fit the "Ideal," "Acceptable," or "Outside" parameters for the role.
    \item \textbf{Promotion Velocity:} Rewards candidates showing rapid progression (e.g., multiple "Senior" or "Lead" titles in a short span).
    \item \textbf{Quality over Quantity:} Focuses on what the candidate has achieved during their tenure, rather than purely counting months.
  \end{itemize}
\end{frame}

% Slide 10
\begin{frame}{Dynamic Career Sequence Alignment}
  \framesubtitle{Applying Bioinformatics to Resumes}
  \begin{itemize}
    \item \textbf{Algorithm:} Adapts the Needleman-Wunsch sequence alignment algorithm.
    \item \textbf{Ideal Sequence:} [Junior $\rightarrow$ Mid $\rightarrow$ Senior $\rightarrow$ Lead $\rightarrow$ Principal]
    \item \textbf{Execution:} Aligns a candidate's actual job titles against the ideal progression.
    \item \textbf{Scoring:} 
    \begin{itemize}
      \item Match Reward (+1.0)
      \item Gap Penalty (-0.5 for career stagnation)
      \item Mismatch Penalty (for erratic jumps)
    \end{itemize}
  \end{itemize}
\end{frame}

% Slide 11
\begin{frame}{Corporate Pedigree via PageRank}
  \begin{itemize}
    \item \textbf{The Concept:} Not all experience is equal. Moving from highly-selective companies indicates a strong pedigree.
    \item \textbf{Implementation:} We build a directed transition graph ($Company_A \rightarrow Company_B$) across the entire candidate pool.
    \item \textbf{PageRank Execution:} Running a localized PageRank determines the "authority" of organizations dynamically, without hardcoding lists.
    \item \textbf{Result:} Candidates coming from high-authority nodes receive a strategic score boost.
  \end{itemize}
\end{frame}

% Slide 12
\begin{frame}{Pillar 3: Behavioral Intent (20\%)}
  \framesubtitle{Measuring Engagement}
  \begin{itemize}
    \item \textbf{Recruiter Response Rate:} Candidates who historically reply to outreach rank higher.
    \item \textbf{Open-to-Work Status:} Active job seekers get a meaningful visibility boost.
    \item \textbf{Platform Activity:} Recent logins and updates signal genuine interest.
    \item \textbf{Notice Period:} Shorter notice periods are rewarded; heavily extended ones receive slight penalties depending on the role urgency.
  \end{itemize}
\end{frame}

% Slide 13
\begin{frame}{Pillar 4: Location Intelligence (10\%)}
  \framesubtitle{Geographical Alignment}
  \begin{itemize}
    \item \textbf{Configurable Zones:} Roles define preferred and acceptable cities.
    \item \textbf{Preferred Cities (e.g., Pune, Delhi NCR):} Grants full location score (1.0).
    \item \textbf{Acceptable Cities (e.g., Bangalore, Mumbai):} Grants moderate score (0.7).
    \item \textbf{Remote / Unlisted:} Grants baseline score (0.4) to avoid entirely eliminating remote-willing candidates.
  \end{itemize}
\end{frame}

% Slide 14
\begin{frame}{System Integrity \& Anti-Gaming}
  \framesubtitle{Defeating the Keyword Stuffers}
  \begin{itemize}
    \item Our system automatically filters deceptive profiles into a "honeypot" quarantine.
    \item \textbf{Impossible Claims:} Flags "Expert" proficiency attached to 0 months of usage.
    \item \textbf{Timeline Inconsistencies:} Flags claims of 10+ years of experience that do not map to the actual career timeline provided.
    \item Protects the integrity of the recruiter's shortlist.
  \end{itemize}
\end{frame}

% Slide 15
\begin{frame}{Shannon Entropy Anomaly Detection}
  \framesubtitle{Statistical Profiling}
  \begin{itemize}
    \item \textbf{The Theory:} Natural human language has a predictable distribution of vocabulary.
    \item \textbf{The Math:} We calculate the Shannon Entropy of the resume text.
    $$H(X) = -\sum P(x_i) \log_2 P(x_i)$$
    \item \textbf{The Trigger:} Profiles whose vocabulary entropy deviates by more than $2\sigma$ from the population mean are automatically flagged.
    \item \textbf{The Result:} Successfully catches machine-generated, keyword-stuffed, and white-fonted resumes.
  \end{itemize}
\end{frame}

% Slide 16
\begin{frame}{High-Level System Architecture}
  \begin{itemize}
    \item \textbf{Two Primary Layers:}
    \begin{enumerate}
      \item Parallelized Python Engine (Offline CLI)
      \item Interactive Recruiter Sandbox (Frontend UI)
    \end{enumerate}
    \item \textbf{Ingestion:} O(N) streaming scanner restricts memory footprint under 200MB.
    \item \textbf{Pre-Scan Phase:} Calculates global TF-IDF weights and entropy baselines.
    \item \textbf{Execution:} Multiprocessing pool shards candidates across CPU cores.
  \end{itemize}
\end{frame}

% Slide 17
\begin{frame}{The Parallelized Scoring Engine}
  \begin{itemize}
    \item \textbf{Speed:} Processes batches of candidates concurrently across all available cores.
    \item \textbf{Benchmark:} 100,000+ candidate profiles ranked in \textbf{under 35 seconds} on a standard CPU.
    \item \textbf{Efficiency:} Achieves execution well within the 5-minute hackathon limit.
    \item \textbf{Deterministic Sort:} Ranks by $(-round(score, 4), candidate\_id)$ guaranteeing identical results on every rerun.
  \end{itemize}
\end{frame}

% Slide 18
\begin{frame}{Explainable AI (XAI)}
  \framesubtitle{Generating Deterministic Rationale}
  \begin{itemize}
    \item \textbf{No Black Boxes:} Every shortlist placement must be justifiable.
    \item \textbf{Context-Aware Generation:} The engine writes human-readable rationale based on actual matched parameters.
    \item \textbf{Example Output:} \textit{"Exceptional technical fit, covering 4/5 core semantic clusters. 6.5 years of steady progression. Background includes Tier-1 product companies."}
    \item Builds trust between the algorithm and the recruiter.
  \end{itemize}
\end{frame}

% Slide 19
\begin{frame}{Fully Offline Operation}
  \framesubtitle{Meeting Strict Privacy and Latency Constraints}
  \begin{itemize}
    \item \textbf{Zero API Calls:} The entire pipeline runs using Python's standard library.
    \item \textbf{No OpenAI/Cloud Dependency:} Eliminates network latency, API rate limits, and token costs.
    \item \textbf{Security:} Enterprise-ready approach where candidate PII never leaves the local machine.
    \item Complies 100\% with sandbox environment restrictions.
  \end{itemize}
\end{frame}

% Slide 20
\begin{frame}{Configurable Job Descriptions}
  \begin{itemize}
    \item \textbf{Adaptability:} No code changes required to hire for a different role.
    \item \textbf{JSON Driven:} All role requirements exist in \texttt{data/jd\_rules.json}.
    \item \textbf{Customization:} Recruiters can easily modify:
    \begin{itemize}
      \item Target Role Title
      \item Experience Minimums and Maximums
      \item Semantic Skill Clusters
      \item Preferred Locations and Company Tiers
    \end{itemize}
  \end{itemize}
\end{frame}

% Slide 21
\begin{frame}{The Recruiter Dashboard UI}
  \framesubtitle{Zero Setup, Maximum Insight}
  \begin{itemize}
    \item \textbf{The Technology:} A pure HTML/JS single-file application. No Node.js, no build steps.
    \item \textbf{Accessibility:} Runs directly in any modern browser immediately after the CLI engine finishes.
    \item \textbf{Purpose:} Translates the raw CSV output into an interactive, visual sandbox for recruiters.
  \end{itemize}
\end{frame}

% Slide 22
\begin{frame}{Dashboard Features in Detail}
  \begin{itemize}
    \item \textbf{CSV Upload:} Drag-and-drop the generated \texttt{team\_submission.csv}.
    \item \textbf{Candidate Cards:} Visual display of rank, score, and the AI rationale.
    \item \textbf{Click-to-Expand:} Opens a detailed scoring breakdown with individual progress bars for Skills, Experience, and Intent.
    \item \textbf{Export Capabilities:} One-click re-export of filtered/sorted results.
  \end{itemize}
\end{frame}

% Slide 23
\begin{frame}{Active Learning Loop (RLRF)}
  \framesubtitle{Continuous Improvement}
  \begin{itemize}
    \item \textbf{Recruiter Feedback:} Dashboard allows recruiters to "Invite" or "Decline" candidates.
    \item \textbf{Weight Adjustment:} The system uses a single-layer perceptron training step to update category weights based on user preference.
    \item \textbf{Adaptation:} Over time, the local weights shift to match the specific hiring manager's implicit preferences.
  \end{itemize}
\end{frame}

% Slide 24
\begin{frame}{Performance Metrics \& Constraints}
  \begin{table}
    \centering
    \begin{tabular}{@{}ll@{}}
      \toprule
      \textbf{Metric} & \textbf{Value} \\
      \midrule
      Dataset size & 100,000+ candidate profiles \\
      Processing time & $\sim$34 seconds \\
      Memory usage & $<$ 500 MB RAM \\
      External API calls & Zero \\
      Dependencies & Pure Python stdlib (Core Engine) \\
      Hackathon Time Limit & 5 minutes (We are 8$\times$ faster) \\
      \bottomrule
    \end{tabular}
  \end{table}
\end{frame}

% Slide 25
\begin{frame}{Conclusion \& Impact}
  \begin{itemize}
    \item \textbf{For Recruiters:} Drastically reduces time-to-shortlist from days to 35 seconds.
    \item \textbf{For Candidates:} Rewards actual career growth, system design skills, and intent---not just resume hacking.
    \item \textbf{For Redrob:} A highly scalable, production-ready heuristic engine adaptable to any JD via a simple configuration.
  \end{itemize}
  \vspace{2em}
  \begin{center}
    \textbf{Thank you! Check out our GitHub Repo.}
  \end{center}
\end{frame}

\end{document}
